\section{Convergence Results and Limit Theorems}

\subsection{Convergence of Random Variable}

\begin{definition}[Convergence in Probability]
    \label{def:convergence-in-probability}%
    A sequence of random variables \(\pqty{X_n}_{n \in \NN}\) \emph{converges} to a random variable \(X\) \emph{in probability}, denoted as
    \[
        X_n \toProb X,
    \]
    if
    \[
        \forall \epsilon > 0: \Prob\pqty{\abs{X_n - X} > \epsilon} \to 0
    \]
    as \(n \to \infty\).
\end{definition}

\begin{theorem}[Weak Law of Large Numbers]
    \label{thm:weak-law-of-large-numbers}%
    Let \(\pqty{X_n}\) be independent and identically distributed distributions with finite mean \(\mu\). Let
    \[
        S_n = X_1 + \cdots + X_n
    \]
    for all \(n \geq 1\). Then we have
    \[
        \frac{S_n}{n} \toProb \mu
    \]
    as \(n \to \infty\). In other words,
    \[
        \Prob\pqty{\abs{\frac{S_n}{n} - \mu} > \epsilon} \to 0
    \]
    as \(n \to \infty\).
\end{theorem}

\begin{proof}[for Finite Variance]
    We give the proof when there is a finite variance. This is not a necessary assumption, but we are using this fact in our proof.

    Let \(\sigma^2\) be the variance and is finite. Then we have
    \begin{align*}
        \Prob\pqty{\abs{\frac{S_n}{n} - \mu} > \epsilon} &= \Prob\pqty{\abs{S_n - n \mu} > n \epsilon} \\
        &\leq \frac{\Var\pqty{S_n}}{n^2 \epsilon^2} \\
        &= \frac{n \sigma^2}{n^2 \epsilon^2} \\
        &= \frac{\sigma^2}{n \epsilon^2} \\
        &\to 0,
    \end{align*}
    by Chebyshev's Inequality, as desired.
\end{proof}

\begin{definition}[Almost Sure Convergence]
    \label{def:almost-sure-convergence}%
    Let \(\pqty{X_n}_{n \in \NN}\) be a sequence of random variables, and let \(X\) be a random variable. We say that \(\pqty{X_n}\) converges to \(X\) \emph{almost surely} (a.s.) or \emph{with probability 1}, if
    \[
        \Prob\pqty{\lim_{n \to \infty} X_n = X} = 1.
    \]

    We write \(X_n \to X\) as \(n \to \infty\) a.s., or \(X_n \toAlmostSure X\).
\end{definition}

\begin{remark}
    The notation
    \[
        \lim_{n \to \infty} X_n = X
    \]
    means we have
    \[
        \forall \epsilon > 0, \exists n_0 \in \NN, \forall n \geq n_0: \abs{X_n - X} \leq \epsilon.
    \]

    When we say this has probability \(1\), this means that there is some set \(A \in \eventSpace\), with \(\Prob\pqty{A} = 1\), such that for all \(\omega \in A\), the above condition holds.
\end{remark}

This notion of convergence is strictly stronger than our previous notion.

\begin{claim}[Almost Sure Convergence to \(0\) implies Convergence in Probability]
    \label{claim:almost-sure-convergence-implies-convergence-in-probability}%
    If \(X_n \toAlmostSure 0\) as \(n \to \infty\), then \(X_n \toProb 0\) as \(n \to \infty\).
\end{claim}

\begin{proof}
    We need to show that
    \[
        \Prob\pqty{\abs{X_n} \leq \epsilon} \to 1.
    \]

    We have
    \[
        \Prob\pqty{\abs{X_n} \leq \epsilon} \geq \Prob\pqty{\bigcap_{m = n}^{\infty} \Bqty{\abs{X_m} \leq \epsilon}} = \Prob\pqty{A_n}.
    \]

    Note that this sequence of events is increasing, i.e., \(A_n \subseteq A_{n + 1}\).

    Therefore,
    \[
        \Prob\pqty{A_n} \to \Prob\pqty{\bigcup_{n \in \NN} A_n} = \Prob\pqty{\bigcup_{n \in \NN} \bigcap_{m = n}^{\infty} \Bqty{\abs{X_m} \leq \epsilon}}.
    \]  

    Therefore,
    \[
        \lim_{n \to \infty} \Prob\pqty{\abs{X_n} \geq \epsilon} \geq \Prob\pqty{\forall \epsilon > 0: \abs{X_m} \leq \epsilon \text{ for all sufficiently large }m} = 1
    \]
    which finishes the proof.
\end{proof}

\begin{theorem}[Strong Law of Large Numbers (SLLN)]
    \label{thm:strong-law-of-large-numbers}%
    Let \(\pqty{X_n}\) be independent and identically distributed distributions with finite mean \(\mu\). Let
    \[
        S_n = X_1 + \cdots + X_n
    \]
    for all \(n \geq 1\). Then we have
    \[
        \frac{S_n}{n} \toAlmostSure \mu
    \]
    as \(n \to \infty\).
\end{theorem}

\begin{proof}[(Non-Examinable)]
    Assume also that the fourth moment is finite, so \(\Expt\bqty{X_1^4} < \infty\). This is not necessary for the theorem, but our proof requires this.

    We set \(Y_i = X_i - \mu\), then \(\Expt\bqty{Y_i} = 0\). Its moment satisfies
    \[
        \Expt\bqty{Y_1^4} = \Expt\bqty{\pqty{X_1 - \mu}^4} \leq 2^4 \cdot \pqty{\Expt\bqty{X_1^4} + \mu^4} < \infty.
    \]

    Let \(S_n = \sum_{i = 1}^{n} Y_i\). We want to show that \(\frac{S_n}{n} \toAlmostSure 0\).

    Our goal to show that
    \[
        \sum_{n = 1}^{\infty} \pqty{\frac{S_n}{n}}^4 < \infty
    \]
    has probability 1. This then implies that \(\frac{S_n}{n} \to 0\) with probability 1.

    It suffices to show that
    \[
        \Expt\bqty{\sum_{n = 1}^{\infty} \pqty{\frac{S_n}{n}}^4} < \infty.
    \]

    We have
    \begin{align*}
        \Expt\bqty{S_n^4} &= \Expt\bqty{\pqty{Y_1 + \cdots + Y_n}^4} \\
        &= \Expt\bqty{\sum_{i = 1}^{n} Y_i^4 + \sum_{1 \leq i < j \leq n} Y_i^2 Y_j^2 + R}
    \end{align*}
    where \(R\) is a sum of terms of the form \(Y_i^3 Y_j, Y_i^2 Y_j Y_k, Y_i Y_j Y_k Y_l\) for distinct indices \(i, j, k, l\). But \(R\) has expectation zero due to independence and zero mean.

    We also have,
    \[
        \Expt\bqty{Y_i^2 Y_j^2} = \pqty{\Expt\bqty{Y_1^2}}^2 \leq \Expt\bqty{Y_i^4} < \infty
    \]
    by Jensen's Inequality. Hence,
    \[
        \Expt\bqty{S_n^4} \leq n \Expt\bqty{Y_i^4} + 6 \cdot \frac{n \pqty{n - 1}}{2} \Expt\bqty{Y_i}^4 \leq 3n^2 \Expt\bqty{Y_i^4}.
    \]

    Hence, we have
    \[
        \Expt\bqty{\sum_{n = 1}^{\infty} \pqty{\frac{S_n}{n}}^4} = \sum_{n = 1}^{\infty} \frac{1}{n^4} \Expt\bqty{S_n^4} \leq \sum_{n = 1}^{\infty} \frac{3 \Expt\bqty{Y_i}^4}{n^2} < \infty
    \]
    which finishes the proof.
\end{proof}

\begin{definition}[Converge in Distribution]
    \label{def:convergence-in-distribution}%    
    A sequence of random variables \(\pqty{X_n}_{n \in \NN}\) \emph{converges} to a random variable \(X\) \emph{in distribution}, denoted as
    \[
        X_n \toDist X,
    \]
    if
    \[
        \Prob\pqty{X_n \leq x} = F_{X_n} \pqty{x} \to F_X \pqty{x} = \Prob\pqty{X \leq x}
    \]
    for all \(x \in \RR\) where \(F_X\) is continuous.
    
    In other words, the cumulative distribution functions of \(X_n\) converges to the cumulative distribution function of \(X\), pointwise, wherever the cumulative distribution function of \(X\) is continuous.
\end{definition}

\begin{claim}[Convergence in Probability implies Convergence in Distribution]
    \label{claim:convergence-in-probability-implies-converges-in-distribution}%
    If \(X_n \toProb X\) as \(n \to \infty\), then \(X_n \toDist X\) as \(n \to \infty\).
\end{claim}

\begin{proof}
    Let \(x\) be a continuity point of \(F_X\). Let \(\epsilon > 0\) be arbitrary. We have
    \[
        \Bqty{X \leq x - \epsilon} \subseteq \Bqty{X_n \leq x} \cup \Bqty{\abs{X_n - X} > \epsilon},
    \]
    and
    \[
        \Bqty{X_n \leq x} \subseteq \Bqty{X \leq x + \epsilon} \cup \Bqty{\abs{X_n - X} > \epsilon}.
    \]

    Hence,
    \[
        \Prob\pqty{X \leq x - \epsilon} \leq \Prob\pqty{X_n \leq x} + \Prob\pqty{\abs{X_n - X} > \epsilon},
    \]
    and taking the limit inferior as \(n \to \infty\) gives
    \[
        F_X \pqty{x - \epsilon} \leq \liminf_{n \to \infty} F_{X_n} \pqty{x}.
    \]

    Similarly,
    \[
        \Prob\pqty{X_n \leq x} \leq \Prob\pqty{X \leq x + \epsilon} = \Prob\pqty{\abs{X_n - X} > \epsilon}
    \]
    and taking the limit superior gives
    \[
        \limsup_{n \to \infty} F_{X_n} \pqty{x} \leq F_X \pqty{x + \epsilon}.
    \]

    Hence,
    \[
        F_{X} \pqty{x - \epsilon} \leq \liminf_{n \to \infty} F_{X_n} \pqty{x} \leq \limsup_{n \to \infty} F_{X_n} \pqty{x} \leq F_X \pqty{x + \epsilon}.
    \]

    Letting \(\epsilon \to 0\), since \(x\) is a continuity point, we have
    \[
        \lim_{\epsilon \to 0} F_X\pqty{x - \epsilon} = \lim_{\epsilon \to 0} F_X\pqty{x + \epsilon} = F_X\pqty{x}
    \]
    and hence
    \[
        \liminf_{n \to \infty} F_{X_n} \pqty{x} = \limsup_{n \to \infty} F_{X_n} \pqty{x} = F_X \pqty{x}
    \]
    as desired.
\end{proof}

\subsection{Central Limit Theorem}

Let \(X_1, X_2, \ldots\) be independent and identically distributed with mean \(\mu\) and variance \(\sigma^2\), both of which are finite. Set \(S_n = X_1 + X_2 + \cdots + X_n\).

By \autoref{thm:strong-law-of-large-numbers} (Strong Law of Large Numbers), we expect that \(S_n \approx n \mu\) for large \(n\). We also have the variance and mean of \(S_n\) satisfies
\[
    \Var\pqty{S_n - n \mu} = n \sigma^2, \quad \Expt\bqty{S_n - n \mu} = 0,
\]
and hence we have
\[
    \frac{S_n - n \mu}{\sqrt{n \sigma^2}} = \frac{\frac{S_n}{n} - \mu}{\sqrt{\Var\pqty{\frac{S_n}{n} - \mu}}}
\]
has expectation \(0\) and variance \(1\).

If \(X_1, \ldots\) are in fact normal distributions \(\Normal\pqty{\mu, \sigma^2}\), then we have
\[
    S_n - n \mu \sim \Normal\pqty{0, n \sigma^2},
\]
and we have
\[
    \frac{S_n - n \mu}{\sqrt{n \sigma^2}} \sim \Normal\pqty{0, 1}
\]
is the standard normal.

In fact, the normal distribution is \emph{universal} as it appears as the limit of the distribution of \(\frac{S_n - n \mu}{\sigma \sqrt{n}}\), no matter what the distribution of the random variables \(X_i\) is.

\begin{theorem}[Central Limit Theorem]
    \label{thm:clt}%
    Let \(X_1, \ldots\) be independent and identically distributed with finite mean \(\mu\) and finite variance \(\sigma^2\). Let \(S_n = X_1 + \cdots + X_n\), then
    \[
        \frac{S_n - n \mu}{\sigma \sqrt{n}} \toDist Z
    \]
    where \(Z\) is the standard normal \(Z \sim \Normal\pqty{0, 1}\).

    In other words, for all \(x \in \RR\), we have
    \[
        \Prob\pqty{\frac{S_n - n \mu}{\sigma \sqrt{n}} \leq x} \to \int_{\minusinfty}^{x}  \frac{e^{- \flatfrac{y^2}{2}}}{\sqrt{2 \pi}} \dd{y} = \sNormCdf\pqty{x}.
    \]
\end{theorem}

We introduce this theorem on moment generating functions for the proof.

\begin{theorem}[Continuity Property for Moment Generating Functions]
    \label{thm:continuity-property-mgf}%
    Suppose \(\pqty{X_n}\) are random variables with
    \[
        m_n \pqty{\theta} = \Expt\bqty{e^{\theta X_n}},
    \]
    for \(\theta \in \RR\), and \(X\) be a random variable with
    \[
        m \pqty{\theta} = \Expt\bqty{e^{\theta X}}.
    \]

    Assume \(m\pqty{\theta}\) is finite for some \(\theta \neq 0\). If \(m_n \pqty{\theta} \to m\pqty{\theta}\) for all \(\theta \in \RR\), then \(X_n \toDist X\) as \(n \to \infty\).
\end{theorem}

\begin{proof}[of \autoref{thm:clt} (Central Limit Theorem)]
    Consider
    \[
        Y_i = \frac{X_i - \mu}{\sigma},
    \]
    then \(\Expt\bqty{Y_i} = 0\) and \(\Var\pqty{Y_i} = 1\). Hence, it suffices to prove the theorem for \(X_1, X_2, \ldots\) with mean \(0\) and variance \(1\). We need to show
    \[
        \frac{S_n}{\sqrt{n}} \toDist \Normal\pqty{0, 1}
    \]
    as \(n \to \infty\).

    Furthermore, we assume there is some \(\delta > 0\), such that
    \[
        \Expt\bqty{e^{\delta X_1}} + \Expt\bqty{e^{- \delta X_1}} < \infty.
    \]

    Let \(m\pqty{\theta} = \Expt\bqty{e^{\theta X_1}}\). By \autoref{thm:continuity-property-mgf} (Continuity Property of Moment Generating Function), it suffices to show
    \[
        \Expt\bqty{e^{\theta \frac{S_n}{\sqrt{n}}}} \to \Expt\bqty{e^{\theta Z}} = e^{\frac{\theta^2}{2}}.
    \]

    Since the \(X_i\)s are independent and identically distributed with moment generating functions \(m\), we have
    \[
        \Expt\bqty{e^{\theta \frac{S_n}{\sqrt{n}}}} = \pqty{m\pqty{\frac{\theta}{\sqrt{n}}}}^{n}.
    \]

    Hence,
    \[
        m\pqty{\frac{\theta}{\sqrt{n}}} = \Expt\bqty{e^{\theta \frac{S_n}{\sqrt{n}}}} = 1 + \frac{\theta^2}{2n} + \Expt\bqty{\sum_{k \geq 3} \frac{\theta^k}{\pqty{\sqrt{n}}^k} \cdot \frac{1}{k!} X_1^k}.
    \]

    We claim that the expectation of the remaining part has
    \[
        \abs{\Expt\bqty{\sum_{k \geq 3} \theta^k \cdot \frac{1}{k!} X_1^k}} = \littleO\pqty{\abs{\theta}^2}
    \]
    as \(\theta \to 0\). If this is indeed true, we have
    \[
        m\pqty{\frac{\theta}{\sqrt{n}}} = 1 + \frac{\theta^2}{2n} + \littleO\pqty{\frac{\theta^2}{n}},
    \]
    and hence
    \[
        \pqty{m \pqty{\frac{\theta}{\sqrt{n}}}}^n \to e^{\flatfrac{\theta^2}{2}}
    \]
    as \(n \to \infty\).

    Indeed, let \(\abs{\theta} \leq \frac{\delta}{2}\), then we have
    \begin{align*}
        \abs{\Expt\bqty{\sum_{k \geq 3} \frac{\theta^k X_1^k}{k!}}} &\leq \Expt\bqty{\sum_{k \geq 3} \frac{\abs{\theta X_1}^k}{k!}} \\
        &= \Expt\bqty{\abs{\theta X_1}^3 \sum_{k \geq 0} \frac{\abs{\theta X_1}^k}{\pqty{k + 3}!}} \\
        &\leq \Expt\bqty{\abs{\theta X_1}^3 \exp\pqty{\abs{\theta X_1}}} \\
        &\leq \Expt\bqty{\abs{\theta X_1}^3 \cdot \exp\pqty{\frac{\delta}{2} \cdot \abs{X_1}}}.
    \end{align*}

    Observe that
    \[
        \abs{\theta X_1}^3 = \abs{\theta}^3 \cdot \frac{\abs{\frac{\delta}{2} X_1}^3}{3!} \cdot \frac{3!}{\pqty{\frac{\delta}{2}}^3} \leq C \cdot \abs{\theta}^3 \cdot e^{\frac{\delta}{2} \abs{X_1}}
    \]
    where
    \[
        C = 3! \cdot \frac{2^3}{\delta^3}
    \]
    is a constant. Hence,
    \begin{align*}
        \Expt\bqty{\abs{\theta X_1}^3 \cdot \exp \pqty{\frac{\delta}{2} \cdot \abs{X_1}}} &\leq C \cdot \abs{\theta}^3 \cdot \Expt\bqty{e^{\delta \abs{X_1}}} \\
        &\leq C \cdot \abs{\theta}^3 \cdot \pqty{\Expt\bqty{e^{\delta X_1} + e^{- \delta X_1}}}
    \end{align*}
    and hence
    \[
        \abs{\Expt\bqty{\sum_{k \geq 3} \frac{\theta^k X_1^k}{k!}}} \leq C' \cdot \abs{\theta}^3
    \]
    for some constant \(C'\).
\end{proof}

A consequence of \autoref{thm:clt} (Central Limit Theorem) is
\[
    S_n \dot{\sim} \Normal\pqty{n \mu, n \sigma^2}
\]
for \(n\) large, and we have
\[
    s_N \approx n \mu + \sigma \sqrt{n} Z
\]
for a standard normal \(Z\), for large \(n\).

\begin{example}[Application of Central Limit Theorem]
    \label{ex:clt}%
    \begin{itemize}
        \item Let \(S_n \sim \Binomial\pqty{n. p}\), then \(S_n = X_1 + \cdots + X_n\) for \(X_i\) being independent and identically distributed \(\Bernoulli \pqty{p}\).

        Hence,
        \[
            \frac{S_n - np}{\sqrt{n p \pqty{1 - p}}} \toDist \Normal\pqty{0, 1}.
        \]

        So
        \[
            S_n \dot{\sim} \Normal\pqty{np, np\pqty{1 - p}}
        \]
        for large \(n\).

        \item Suppose instead that \(S_n \sim \Binomial\pqty{n, \frac{\lambda}{n}}\) for \(\lambda > 0\), and let \(S \sim \Poisson\pqty{\lambda}\). Then we have
        \[
            \Prob\pqty{S_n = x} \to \Prob\pqty{S = x}.
        \]

        \item Suppose now that \(S_n \sim \Poisson\pqty{n}\), then \(S_n = X_1 + \cdots + X_n\) for \(X_i\) being independent and identically distributed \(\Poisson \pqty{1}\). Hence,
        \[
            \frac{S_n - n}{\sqrt{n}} \toDist \Normal\pqty{0, 1},
        \]
        and hence
        \[
            S_n \dot{\sim} \Normal\pqty{n, n}.
        \]
    \end{itemize}
\end{example}

\begin{example}[Sampling Error Estimation using Central Limit Theorem]
    Suppose we hold a referendum, and a proportion \(p\) of individuals in inclined to vote `Yes'. In other words, every individual votes `Yes' with probability \(p\), and `No' with probability \(\pqty{1 - p}\), independently of each other.

    We want to estimate \(p\). Sample \(N\) individuals at random, where \(N\) is large, and we record their votes. Let \(X_i\) be \(1\) if the \(i\)th individual voted `Yes', and \(0\) otherwise.

    Set \(S_N = X_1 + \cdots + X_n\) be the number of `Yes' votes, and denote
    \[
        \hat{P}_N = \frac{S_N}{N} \toAlmostSure p
    \]
    by \autoref{thm:strong-law-of-large-numbers} (Strong Law of Large Numbers).
    
    We have \(S_N \sim \Binomial\pqty{N, p}\).
    
    Now, our goal is to estimate \(p\) with an accuracy of \(\pm 4\%\) with probability \(\geq 0.99\). How large should \(N\) to be to ensure this?

    By \autoref{thm:clt} (Central Limit Theorem),
    \[
        S_N \approx Np + \sqrt{N p \pqty{1 - p}} Z
    \]
    for \(Z \sim \Normal\pqty{0, 1}\), where \(N\) is large. Rephrasing the question, we want to find \(N\), such that
    \[
        \Prob\pqty{\abs{\hat{P}_N - p} \geq 0.04} \leq 0.01.
    \]

    Hence, we have
    \[
        \hat{P}_N = \frac{S_N}{N} \approx p + \frac{\sqrt{p \pqty{1 - p}}}{\sqrt{N}} \cdot Z
    \]  
    for large \(N\). So we have
    \[
        \abs{\hat{P}_N} - p \approx \frac{\sqrt{p \pqty{1 - p}}}{\sqrt{N}} \cdot \abs{Z}.
    \]

    Furthermore, we have
    \[
        \Prob\pqty{\abs{Z} \geq z} = 2 \pqty{1 - \sNormCdf\pqty{z}}.
    \]

    Looking at standard tables, we have for \(z = 2.58\), we have
    \[
        \Prob\pqty{\abs{Z} \geq z} = 0.01.
    \]

    Hence, we may choose a sufficiently large \(N\), such that
    \[
        \frac{\sqrt{N}}{\sqrt{p \pqty{1 - p}}} \cdot 0.04 \geq 2.58.
    \]

    But we want this to hold for any \(p\), and note that
    \[
        \sqrt{p \pqty{1 - p}} \leq \frac{1}{2},
    \]
    so taking \(N \approx 1040\) should be sufficient.
\end{example}